% Standalone problem document, generated from the adjacent README.md.
% Compile with XeLaTeX. Mathematical content is maintained in Markdown.
\documentclass[11pt,a4paper]{article}
\usepackage[fontsize=10.5pt]{scrextend}
\usepackage[margin=22mm,headheight=15pt,headsep=9mm,footskip=11mm]{geometry}
\usepackage{fontspec}
\setmainfont{texgyrepagella}[Extension=.otf,UprightFont=*-regular,BoldFont=*-bold,ItalicFont=*-italic,BoldItalicFont=*-bolditalic]
\setsansfont{texgyreheros}[Extension=.otf,UprightFont=*-regular,BoldFont=*-bold,ItalicFont=*-italic,BoldItalicFont=*-bolditalic]
\setmonofont{texgyrecursor}[Extension=.otf,UprightFont=*-regular,BoldFont=*-bold,ItalicFont=*-italic,BoldItalicFont=*-bolditalic,Scale=MatchLowercase]
\usepackage{amsmath,amssymb}
\usepackage{unicode-math}
\setmathfont{texgyrepagella-math.otf}
\usepackage{xcolor,microtype,enumitem,needspace,fancyhdr}
\usepackage{bookmark}
\definecolor{ink}{HTML}{17344B}
\definecolor{accent}{HTML}{197D80}
\definecolor{muted}{HTML}{596773}
\hypersetup{colorlinks=true,linkcolor=accent,urlcolor=accent,pdftitle={MD-03 - The
Komlós discrepancy
conjecture},pdfauthor={Open Problems in Numerical Linear Algebra}}
\urlstyle{same}
\setlength{\parindent}{0pt}
\setlength{\parskip}{5pt plus 1pt minus 1pt}
\linespread{1.04}
\setlength{\emergencystretch}{3em}
\widowpenalty=10000
\clubpenalty=10000
\displaywidowpenalty=10000
\allowdisplaybreaks[1]
\setlist{nosep,leftmargin=1.5em}
\providecommand{\tightlist}{\setlength{\itemsep}{0pt}\setlength{\parskip}{0pt}}
\providecommand{\pandocbounded}[1]{#1}
\pagestyle{fancy}
\fancyhf{}
\fancyhead[L]{\sffamily\footnotesize\color{muted}OPEN PROBLEMS IN NUMERICAL LINEAR ALGEBRA}
\fancyhead[R]{\sffamily\bfseries\footnotesize\color{accent}MD-03}
\fancyfoot[L]{\parbox[b]{0.9\textwidth}{\sffamily\fontsize{8}{10}\selectfont\color{muted}Editorial ratings; a bounded search cannot certify that a problem remains unsolved.\\Literature check: 2026-09-11}}
\fancyfoot[R]{\sffamily\footnotesize\color{muted}\thepage}
\renewcommand{\headrulewidth}{0.3pt}
\renewcommand{\headrule}{\hbox to\headwidth{\color{accent}\leaders\hrule height \headrulewidth\hfill}}
\makeatletter
\renewcommand{\section}{\Needspace{5\baselineskip}\@startsection{section}{1}{0pt}{12pt plus 2pt minus 2pt}{4pt}{\sffamily\bfseries\large\color{ink}}}
\renewcommand{\subsection}{\Needspace{4\baselineskip}\@startsection{subsection}{2}{0pt}{9pt plus 1pt minus 1pt}{3pt}{\sffamily\bfseries\normalsize\color{ink}}}
\makeatother
\setcounter{secnumdepth}{0}
\begin{document}
{\sffamily\small\color{muted}Matrix discrepancy and optimization\par}
\vspace{3pt}
{\raggedright\hyphenpenalty=10000\exhyphenpenalty=10000\sffamily\bfseries\fontsize{21}{24}\selectfont\color{ink}The
Komlós discrepancy conjecture\par}
\vspace{7pt}
{\sffamily\small\textbf{Difficulty:} extreme\quad\textbf{Importance:} broadly
interesting\par}
{\sffamily\small\bfseries\color{accent}Status: Solved\par}
\vspace{4pt}
{\color{accent}\hrule height 0.8pt}
\vspace{6pt}
\textbf{Provenance:} source-stated conjecture.

\textbf{Rating rationale:} Dimension-free vector balancing is a
foundational discrepancy barrier with wide implications for rounding,
optimization and algorithms.

\subsection{Resolution --- 2026-09-11}\label{resolution-2026-09-11}

\textbf{Affirmative resolution by Shengtao Guo, Ethan X. Fang and Junwei
Lu}, \href{https://arxiv.org/abs/2609.11189v1}{\emph{Vector Balancing
via Directional Total Variation}, arXiv:2609.11189v1}, submitted 10
September 2026, \textbf{Theorem 1.1, p.~1}. It proves the displayed
target with the universal constant \(C=3\sqrt{2\pi}\); its strict bound
is stronger than the requested non-strict inequality. Each column of
\(A\) is one of the theorem's vectors. All positive dimensions and every
column are covered; no algorithmic assertion is required.

\textbf{Application note by George Stepaniants}, Department of Computing
and Mathematical Sciences, California Institute of Technology:
\href{https://github.com/ajt60gaibb/OpenProblemsInNLA/blob/main/matrix-discrepancy-and-optimization/MD-03/solution.md}{proof
by reference}, \textbf{Theorem 1}
(\href{https://github.com/ajt60gaibb/OpenProblemsInNLA/blob/main/references/stepaniants-2026-09-11/manuscripts/md03_md04_proofs.pdf}{PDF}
·
\href{https://github.com/ajt60gaibb/OpenProblemsInNLA/blob/main/references/stepaniants-2026-09-11/manuscripts/md03_md04_proofs.tex}{LaTeX}).
Stepaniants authors this explanatory note; the discrepancy theorem is
credited to Guo, Fang and Lu.

A separate Codex agent checked both the application and the substantive
source proof, including its published eigenvalue-convexity input, and
returned \textbf{PASS}:
\href{https://github.com/ajt60gaibb/OpenProblemsInNLA/blob/main/references/stepaniants-2026-09-11/verification/reviews/MD-03-MD-04-review.md}{independent
review}. This is independent automated-agent verification under the
repository's status definition, not external human peer review or formal
certification. The source remains a preprint and discloses Odin AI use.
The application draft was prepared with ChatGPT. The original statement
and dated history below are retained; the ratings above are historical.
\href{https://github.com/ajt60gaibb/OpenProblemsInNLA/blob/main/references/stepaniants-2026-09-11/README.md}{Submission
record}.

\subsection{Problem statement}\label{problem-statement}

Does a finite constant \(C>0\) exist such that, for every pair of
positive integers \(m,n\) and every matrix
\(A=(a_{ij})\in\mathbb R^{m\times n}\) satisfying \[
\sum_{i=1}^m a_{ij}^2\leq1\qquad(1\leq j\leq n),
\] there is a vector \(x\in\{-1,1\}^n\) for which \[
\|Ax\|_\infty=\max_{1\leq i\leq m}\left|\sum_{j=1}^n a_{ij}x_j\right|\leq C?
\] The same constant must work for every dimension and matrix. The
question asks for existence of a signing; it does not additionally
require an algorithm.

This is a matrix balancing problem: choose column signs so that their
sum is small in every coordinate. Its norm constraint and simultaneous
coordinate control connect discrepancy theory with rounding and linear
optimization.

\subsection{References}\label{references}

\begin{enumerate}
\def\labelenumi{\arabic{enumi}.}
\tightlist
\item
  N. Bansal and H. Jiang, \emph{An Exposition of the
  \(\widetilde O(\log^{1/4}n)\) Bound for the Komlós Problem},
  arXiv:2608.28452 (2026), abstract and introduction, including the
  explicit conjecture and the new dimension-dependent bound.
  \href{https://arxiv.org/html/2608.28452v1}{Paper}.
\item
  N. Bansal and H. Jiang, \emph{Decoupling via Affine
  Spectral-Independence: Beck-Fiala and Komlós Bounds Beyond
  Banaszczyk}, arXiv:2508.03961v2 (2025), abstract and introductory
  results. \href{https://arxiv.org/abs/2508.03961}{Paper}.
\item
  W. Banaszczyk, \emph{Balancing vectors and Gaussian measures of
  n-dimensional convex bodies}, Random Structures \& Algorithms 12(4)
  (1998), pp.~351--360, the main vector-balancing theorem underlying the
  earlier logarithmic bound.
  \href{https://onlinelibrary.wiley.com/doi/abs/10.1002/\%28SICI\%291098-2418\%28199807\%2912\%3A4\%3C351\%3A\%3AAID-RSA3\%3E3.0.CO\%3B2-S}{Article}.
\end{enumerate}

\subsection{Earlier status check ---
2026-09-10}\label{earlier-status-check-2026-09-10}

checked the current arXiv records (2608.28452v1, 2026-08-28;
2508.03961v2, 2025-09-09) and searched ``Komlos conjecture solved proof
September 2026'' and ``Komlos Bansal Jiang 2026 bound''. The new bound
is \(O((\log n)^{1/4}(\log\log n)^{7/4})\) in its asymptotic range, not
a universal constant. No full proof, counterexample, or withdrawal of
these source papers was found. The disproved Hajela conjecture and the
separate matrix Spencer conjecture are not this statement.

\textbf{Audit update (2026-09-10):} Rechecked the August 2026 exposition
and searched for later Komlós resolutions. Its improved
dimension-dependent bound still leaves the universal-constant target
open. This is a bounded literature check, not a proof that no solution
exists.
\end{document}
